\section{Where the estimate enters spectral-gap arguments}\label{sec:spectral}

This section identifies the application of Theorem~\ref{thm:mixed} and explains the surrounding mechanism. The graph-level path enumeration and tangle estimates are separate ingredients from the papers cited below.

\subsection{Random lifts: one permutation per base edge}

In Bordenave's model, choose one orientation of each undirected edge of a finite base graph. Attach an independent uniform permutation of $[n]$ to each chosen edge. Its reverse orientation carries the inverse permutation. A lift edge between fibers is represented by a matrix entry of this permutation, and reversing the edge transposes the same entry.

Orient every factor in a requested product according to the chosen base orientations. Base edges become colors in \eqref{eq:product}. Distinct requested lift edges become distinct cells. Their multiplicities and consistency conditions are unchanged. The parameter $a_1$ of \cite[Proposition 28]{Bordenave} is exactly $h$ here, and its inconsistent centered-singleton count is $b$.

If the original product has $k$ factors, then $a\le k\le\sqrt n$. For $a\ge2$, Theorem~\ref{thm:mixed} gives
\[
|\E P|\le\sqrt e\,2^bn^{-a}(2\sqrt{a/n})^{a_1}
\le\sqrt e\,2^bn^{-a}(3k/\sqrt n)^{a_1}.
\]
The case $a=1$ is covered directly by the theorem. This proves the stated product estimate in Proposition 28, with a stronger consistent-singleton factor.

\subsection{Biregular graphs: the size is the half-edge count}

Take a bipartite configuration model with $n_L$ left vertices of degree $d_L$ and $n_R$ right vertices of degree $d_R$. The two half-edge sets both have size
\[
L=n_Ld_L=n_Rd_R.
\]
A uniform matching between them is a uniform permutation of $[L]$. Its symmetric half-edge matching matrix has the block form
\[
\begin{pmatrix}0&U\\U^{\mathsf T}&0\end{pmatrix},
\]
where $U$ is a uniform $L$ by $L$ permutation matrix. Same-side entries of the matching matrix and its centered version vanish, so any product containing them is zero. Orient each remaining requested matching factor from left to right. Theorem~\ref{thm:mixed} applies with $n=L$, exactly as needed in \cite[preprint Lemma 7, published Lemma 5.6]{BDH}. The centering is $1/L$, not $1/n_L$ or $1/n_R$.

This model differs from an unrestricted perfect matching of a single set of half-edges. That latter model is treated by Bordenave's Proposition 11. The off-diagonal bipartite block is what makes Proposition 28 the correct input here.

\subsection{Why isolated centered factors help a trace expansion}

The non-backtracking matrix is indexed by oriented graph edges. Its entry from $u\to v$ to $x\to y$ is one exactly when $v=x$ and $y\ne u$. Thus powers of this matrix count walks that never immediately reverse an edge.

The matrix is generally not normal. A safe norm estimate uses
\[
\|C\|^{2s}\le\tr\bigl((CC^*)^s\bigr).
\]
Expanding the right side for a centered path matrix produces products of matching entries. A product that uses $a$ distinct matches starts at scale $L^{-a}$. Each consistent centered singleton gives an additional factor of order $\sqrt{a/L}$ in our bound. Repeated cells and conflicting singletons have different costs and must be counted separately.

Choosing locations for a circuit contributes positive powers of the graph size. Matching probabilities cancel most of them. Excess cycles and isolated centered factors supply the savings that remain. The combinatorial argument must sum these estimates over all allowed circuit shapes. The mixed-product theorem handles each fixed pattern, but does not perform this sum.

Short neighborhoods are typically tree-like, with a controlled allowance for a single cycle. The papers restrict to such paths and estimate the remainder from more complicated local tangles. Non-backtracking walks make this separation useful because immediate reversals would create many large, uninformative contributions even on a tree.

\subsection{The algebraic bridge back to adjacency}

For a finite simple graph with adjacency matrix $A$, degree matrix $D$ and non-backtracking matrix $B$, the determinant identity of Ihara and Bass in the cited papers reads
\begin{equation}\label{eq:ihara}
\det(I-uB)=(1-u^2)^{|E|-|V|}\det\bigl(I-uA+u^2(D-I)\bigr).
\end{equation}
For a $d$-regular graph this reduces, away from the exceptional roots, to
\[
\lambda=z+\frac{d-1}{z}.
\]
The circle $|z|=\sqrt{d-1}$ corresponds to the adjacency interval $[-2\sqrt{d-1},2\sqrt{d-1}]$. An adjacency eigenvalue outside that interval produces a non-backtracking root outside the circle.

For a bipartite biregular graph write
\[
A=\begin{pmatrix}0&X\\X^{\mathsf T}&0\end{pmatrix},\qquad
\alpha=d_L-1,\quad\gamma=d_R-1,
\]
and assume $d_L,d_R\ge2$. If $s>0$ is a singular value of $X$, the relevant block of \eqref{eq:ihara} yields
\begin{equation}\label{eq:biregularbridge}
(z^2+\alpha)(z^2+\gamma)=s^2z^2.
\end{equation}
Putting $x=z^2$ gives the scalar map
\begin{equation}\label{eq:scalarbridge}
s^2=x+\frac{\alpha\gamma}{x}+\alpha+\gamma.
\end{equation}
The circle $|x|=\sqrt{\alpha\gamma}$ corresponds to
\[
|\sqrt\alpha-\sqrt\gamma|\le s\le\sqrt\alpha+\sqrt\gamma.
\]
If $s>\sqrt\alpha+\sqrt\gamma$, the quadratic in $x$ has a positive root greater than $\sqrt{\alpha\gamma}$, so $|z|>(\alpha\gamma)^{1/4}$. This is the precise route by which a non-backtracking bound excludes adjacency outliers. Rectangular nullspaces can also produce zero adjacency eigenvalues, which are outside this positive-singular-value calculation.

For example, degrees $(3,7)$ give the positive bulk interval $[\sqrt6-\sqrt2,\sqrt6+\sqrt2]$, approximately $[1.035,3.864]$, while the trivial singular value is $\sqrt{21}\approx4.583$. The spectral theorem controls the remaining singular values asymptotically with high probability. Our proof supplies its matching-product input, not a replacement for all its graph-level estimates.
